%%
%% This is file `sample-sigconf-authordraft.tex',
%% generated with the docstrip utility.
%%
%% The original source files were:
%%
%% samples.dtx  (with options: `all,proceedings,bibtex,authordraft')
%% 
%% IMPORTANT NOTICE:
%% 
%% For the copyright see the source file.
%% 
%% Any modified versions of this file must be renamed
%% with new filenames distinct from sample-sigconf-authordraft.tex.
%% 
%% For distribution of the original source see the terms
%% for copying and modification in the file samples.dtx.
%% 
%% This generated file may be distributed as long as the
%% original source files, as listed above, are part of the
%% same distribution. (The sources need not necessarily be
%% in the same archive or directory.)
%%
%%
%% Commands for TeXCount
%TC:macro \cite [option:text,text]
%TC:macro \citep [option:text,text]
%TC:macro \citet [option:text,text]
%TC:envir table 0 1
%TC:envir table* 0 1
%TC:envir tabular [ignore] word
%TC:envir displaymath 0 word
%TC:envir math 0 word
%TC:envir comment 0 0
%%
%% The first command in your LaTeX source must be the \documentclass
%% command.
%%
%% For submission and review of your manuscript please change the
%% command to \documentclass[manuscript, screen, review]{acmart}.
%%
%% When submitting camera ready or to TAPS, please change the command
%% to \documentclass[sigconf]{acmart} or whichever template is required
%% for your publication.
%%
%%
\documentclass[sigconf, screen, review]{acmart}
%%
%% \BibTeX command to typeset BibTeX logo in the docs
\AtBeginDocument{%
  \providecommand\BibTeX{{%
    Bib\TeX}}}

%% Rights management information.  This information is sent to you
%% when you complete the rights form.  These commands have SAMPLE
%% values in them; it is your responsibility as an author to replace
%% the commands and values with those provided to you when you
%% complete the rights form.
% \setcopyright{acmlicensed}
% \copyrightyear{2018}
\acmYear{2026}
% \acmDOI{XXXXXXX.XXXXXXX}
%% These commands are for a PROCEEDINGS abstract or paper.
\acmConference[ACM MM '26]{the 34th ACM International Conference on Multimedia}
{November 10--14, 2026}{Rio de Janeiro, Brazil}
%%
%%  Uncomment \acmBooktitle if the title of the proceedings is different
%%  from ``Proceedings of ...''!
%%
%%\acmBooktitle{Woodstock '18: ACM Symposium on Neural Gaze Detection,
%%  June 03--05, 2018, Woodstock, NY}
% \acmISBN{978-1-4503-XXXX-X/2018/06}


%%
%% Submission ID.
%% Use this when submitting an article to a sponsored event. You'll
%% receive a unique submission ID from the organizers
%% of the event, and this ID should be used as the parameter to this command.
%%\acmSubmissionID{123-A56-BU3}

%%
%% For managing citations, it is recommended to use bibliography
%% files in BibTeX format.
%%
%% You can then either use BibTeX with the ACM-Reference-Format style,
%% or BibLaTeX with the acmnumeric or acmauthoryear sytles, that include
%% support for advanced citation of software artefact from the
%% biblatex-software package, also separately available on CTAN.
%%
%% Look at the sample-*-biblatex.tex files for templates showcasing
%% the biblatex styles.
%%

%%
%% The majority of ACM publications use numbered citations and
%% references.  The command \citestyle{authoryear} switches to the
%% "author year" style.
%%
%% If you are preparing content for an event
%% sponsored by ACM SIGGRAPH, you must use the "author year" style of
%% citations and references.
%% Uncommenting
%% the next command will enable that style.
%%\citestyle{acmauthoryear}

\usepackage{subcaption}
\usepackage{bm}
\usepackage{multirow}
\usepackage{pifont}
\newcommand{\cmark}{\ding{51}}
\usepackage{booktabs, tabularx}

%%
%% end of the preamble, start of the body of the document source.
\begin{document}

%%
%% The "title" command has an optional parameter,
%% allowing the author to define a "short title" to be used in page headers.
\title{ChainBench-ADD: A Delivery-Aware Dataset and Benchmark for Audio Deepfake Detection}

%%
%% The "author" command and its associated commands are used to define
%% the authors and their affiliations.
%% Of note is the shared affiliation of the first two authors, and the
%% "authornote" and "authornotemark" commands
%% used to denote shared contribution to the research.
\author{Ruiming Wang}
\email{wangrm23@m.fudan.edu.cn}
\affiliation{%
  \institution{Fudan University}
  \city{Shanghai}
  \country{China}
}

\author{Yiliao Song}
\email{lia.song@adelaide.edu.au}
\affiliation{%
  \institution{Adelaide University}
  \city{Adelaide}
  \country{Australia}
}

\author{Ting Dang}
\email{Ting.Dang@unimelb.edu.au}
\affiliation{%
  \institution{The University of Melbourne}
  \city{Melbourne}
  \country{Australia}
}

\author{Bin Wang}
\email{wangbin@nufe.edu.cn}
\affiliation{%
  \institution{Nanjing University of Finance \& Economics}
  \city{Nanjing}
  \country{China}
}

\author{Xin Yu}
\email{xin.yu@adelaide.edu.au}
\affiliation{%
  \institution{Adelaide University}
  \city{Adelaide}
  \country{Australia}
}

%%
%% By default, the full list of authors will be used in the page
%% headers. Often, this list is too long, and will overlap
%% other information printed in the page headers. This command allows
%% the author to define a more concise list
%% of authors' names for this purpose.
% \renewcommand{\shortauthors}{Trovato et al.}

%%
%% The abstract is a short summary of the work to be presented in the
%% article.
\begin{abstract}
Existing audio deepfake datasets have greatly expanded evaluation across generators, languages, and domains, but most remain generation-centric and provide limited support for studying post-generation delivery. In realistic misuse scenarios, forged audio is often altered by platform re-encoding, telephony transmission, or replay-like recapture before detection. Existing delivery-related studies usually model such factors as isolated distortions, rather than structured and ordered delivery routes, and often lack control over transcript and speaker conditions. To address this gap, we introduce ChainBench-ADD, a delivery-aware dataset and benchmark for audio deepfake detection. It treats post-generation delivery itself as a structured benchmark dimension, jointly providing multi-family delivery, explicit ordered chain modeling, and matched-parent controlled evaluation within a unified framework. It organizes delivery into five realistic families using reusable operators, ordered templates, and realized chains, while linking each delivered sample to a clean bona fide or spoof parent under fixed transcript and speaker conditions. The current release contains 941,201 waveforms derived from 55,813 parents and supports five protocol-aware tasks. Experiments show that ChainBench-ADD reveals family-level difficulty differences, sensitivity to controlled local edits, and robustness decay along accumulated delivery paths that conventional pooled evaluation cannot capture. The website is at https://wtalioy.github.io/ChainBench-ADD/.
\end{abstract}

%%
%% The code below is generated by the tool at http://dl.acm.org/ccs.cfm.
%% Please copy and paste the code instead of the example below.
%%
\begin{CCSXML}
<ccs2012>
   <concept>
       <concept_id>10002978.10002997</concept_id>
       <concept_desc>Security and privacy~Intrusion/anomaly detection and malware mitigation</concept_desc>
       <concept_significance>500</concept_significance>
       </concept>
   <concept>
       <concept_id>10002951.10003227.10003251</concept_id>
       <concept_desc>Information systems~Multimedia information systems</concept_desc>
       <concept_significance>500</concept_significance>
       </concept>
 </ccs2012>
\end{CCSXML}

\ccsdesc[500]{Security and privacy~Intrusion/anomaly detection and malware mitigation}
\ccsdesc[500]{Information systems~Multimedia information systems}

%%
%% Keywords. The author(s) should pick words that accurately describe
%% the work being presented. Separate the keywords with commas.
\keywords{Audio Deepfake Detection, Speech Spoofing, Robustness Evaluation, Delivery Chain Modeling, Telephony and Platform Distortions, Replay Simulation}
%% A "teaser" image appears between the author and affiliation
%% information and the body of the document, and typically spans the
%% page.

%%
%% This command processes the author and affiliation and title
%% information and builds the first part of the formatted document.
\begin{teaserfigure}
  \centering
  \includegraphics[width=\textwidth]{figures/overview_figure.png}
  \caption{ChainBench-ADD models post-generation delivery as ordered chains under matched-parent control, going beyond generation-centric benchmarks and isolated delivery effects.}
  \label{fig:teaser}
\end{teaserfigure}

\maketitle

\section{Introduction}
\label{sec:introduction}

Recent advances in text-to-speech, voice cloning, and speech generation have made audio deepfakes increasingly realistic and accessible, turning audio deepfake detection into an important problem in multimedia security. In response, the community has developed a broad range of benchmarks and datasets that substantially improve evaluation across generators, tasks, languages, and domains~\cite{liu2023asvspoof,yi2024add,frank2021wavefake,khalid2021fakeavceleb,muller2022generalize,zhang202partial,sun2023ai,muller2024mlaad,ma2024cfad,wang2026asvspoof,xie2025codecfake,bhagtani2025diffssd,ciobanu2026xmad,wang2025audeter}. These resources have considerably strengthened evaluation on the generation side of audio deepfakes.

However, realistic misuse rarely ends at synthesis. In practice, a forged clip may be uploaded to an online platform, transmitted through a telephony channel, or played aloud and re-recorded before it reaches a listener or an automated system. The audio encountered at deployment is therefore often not pristine generated speech, but speech after post-generation delivery. This creates a mismatch between common benchmark conditions and realistic use scenarios. Some recent studies have started to consider communication effects, codec distortions, or replay-like conditions~\cite{shi2025benchmarking,muller2025replay,delgado2026presentation}, but they mostly examine delivery through isolated effects. While useful, such settings do not fully reflect the fact that post-generation delivery is often a structured process in which multiple operations are combined, parameterized, and ordered. Platform re-encoding, telephony transmission, and replay-like recapture may affect forensic cues differently, and their impact may depend on the particular delivery route rather than on any single operation alone.

These considerations suggest the need for a benchmark that treats post-generation delivery as a structured evaluation dimension: one that spans multiple realistic delivery families, represents delivery as ordered chains rather than isolated perturbations, and supports controlled analysis under matched conditions. Such control is important because, without matched clean parents, delivery effects can be entangled with differences in linguistic content, speaker identity, or source audio quality, making it difficult to attribute detector behavior specifically to delivery.

To this end, we introduce ChainBench-ADD, a delivery-aware dataset and benchmark for audio deepfake detection. It models post-generation delivery through reusable operators, ordered templates, and realized chains across five delivery families: direct, platform-like, telephony, simulated replay, and hybrid. Each delivered sample remains linked to a clean bona fide or spoof parent under matched-parent control, enabling controlled comparison across delivery scenarios. Figure~\ref{fig:teaser} summarizes the benchmark design.

The current release contains over 941k waveforms derived from 55,813 parents. From the shared metadata, we define five evaluation tasks: in-chain detection, three matched local interventions (operator substitution, parameter perturbation, and order swap), and lineage-based delivery robustness. A leave-one-template-out protocol further tests transfer to unseen templates within a family. Together, these components support more realistic and controlled delivery-aware evaluation for audio deepfake detection.

The main contributions of this work are as follows:
\begin{itemize}
    \item We introduce ChainBench-ADD, a large-scale delivery-aware dataset and benchmark for audio deepfake detection that models post-generation delivery as structured chains across five realistic families through reusable operators, ordered templates, and realized chains.
    \item We design a matched-parent, metadata-driven benchmark framework that isolates delivery effects under fixed transcript and speaker conditions and supports controlled, protocol-aware evaluation within a unified metadata schema.
    \item We release a protocol-aware resource with explicit lineage annotations, benchmark statistics, and consistency validation, and show through baseline studies that it exposes family-level, local, and path-wise robustness patterns that conventional pooled evaluation cannot capture.
\end{itemize}
\section{Related Work}
\label{sec:related_work}

Recent audio deepfake resources have developed along four complementary directions: challenge benchmarks, diversity-oriented datasets, delivery-oriented robustness studies, and reproducibility frameworks. Rather than comparing only dataset scale, Table~\ref{tab:rw_compare} focuses on the aspects most relevant to this work: delivery scope, structured composition, and matched-parent control.

\begin{table*}[t]
\centering
\footnotesize
\setlength{\tabcolsep}{3pt}
\renewcommand{\arraystretch}{1.01}
% \caption{Representative prior resources versus ChainBench-ADD. \cmark denotes explicit support, (\cmark) partial or special-case support, \textemdash\ absence, and \textsc{n/a} not applicable because the work is primarily an evaluation framework.}
\caption{Representative prior resources versus ChainBench-ADD. \cmark\ denotes explicit support, (\cmark) partial support, \textemdash\ absence, and \textsc{n/a} not applicable because the work is primarily an evaluation framework.}
\label{tab:rw_compare}
\vspace{-0.5em}
\begin{tabular}{@{}p{2.95cm}p{2.45cm}p{4.85cm}cc@{}}
\toprule
\textbf{Resource} & \textbf{Focus} & \textbf{Post-generation delivery} & \textbf{Ordered comp.} & \textbf{Matched-parent control} \\
\midrule
ASVspoof 2021~\cite{liu2023asvspoof} & challenge benchmark & codec/channel/room effects in challenge tracks & \textemdash & \textemdash \\
ADD 2023~\cite{yi2024add} & task benchmark & not delivery-centric & \textemdash & \textemdash \\
CFAD~\cite{ma2024cfad} & generator/domain range & noise + codec condition sets & \textemdash & \textemdash \\
CodecFake~\cite{xie2025codecfake} & ALM/codec coverage & generator-side codec effects & \textemdash & \textemdash \\
ADD-C~\cite{shi2025benchmarking} & comms robustness & codec + packet-loss conditions & (\cmark) & \textemdash \\
ReplayDF~\cite{muller2025replay} & replay robustness & playback + re-recording pipeline & (\cmark) & \textemdash \\
AUDDT~\cite{zhu2025auddt} & evaluation toolkit & \textsc{n/a} & \textsc{n/a} & \textsc{n/a} \\
Spoof-SUPERB~\cite{ali2026superb} & SSL benchmark & \textsc{n/a} & \textsc{n/a} & \textsc{n/a} \\
\textbf{ChainBench-ADD (ours)} & \textbf{structured delivery} & \textbf{multi-family post-generation delivery} & \textbf{\cmark} & \textbf{\cmark} \\
\bottomrule
\end{tabular}
\end{table*}

Early evaluation standards were shaped largely by the ASVspoof series. ASVspoof 2021~\cite{liu2023asvspoof} formalized logical-access, physical-access, and deepfake evaluation while incorporating compression, transmission, and acoustic propagation effects; ASVspoof 5~\cite{wang2026asvspoof} further broadened speaker, acoustic, and attack diversity. In parallel, ADD 2023~\cite{yi2024add} expanded evaluation beyond utterance-level bona fide/spoof discrimination to include manipulation localization and algorithm recognition, and PartialSpoof~\cite{zhang202partial} focused on partially manipulated utterances. Earlier resources such as WaveFake~\cite{frank2021wavefake} and FakeAVCeleb~\cite{khalid2021fakeavceleb} also helped establish benchmarked evaluation for synthetic and multimodal deepfakes.

A second line of work emphasizes broader generator, language, and domain coverage. In-the-Wild~\cite{muller2022generalize} highlighted the gap between laboratory benchmarks and real online media. MLAAD~\cite{muller2024mlaad} expanded multilingual and multi-generator coverage, LibriSeVoc~\cite{sun2023ai} emphasized vocoder artifacts and self-vocoding conditions, CFAD~\cite{ma2024cfad} introduced a Chinese benchmark under noisy and transcoding conditions, and CodecFake~\cite{xie2025codecfake} focused on codec-based and neural-codec generation. More recent datasets~\cite{huang-etal-2025-speechfake,mawalim2025jmad,ciobanu2026xmad,wang2025audeter} continue this trend toward larger-scale multilingual, cross-domain, and open-world evaluation. Despite that, the main axis of variation in these resources remains the generator, language, or source domain.

More closely related to our setting are delivery-oriented robustness studies. ADD-C~\cite{shi2025benchmarking} shows that codec compression, packet loss, and communication-channel effects can noticeably affect detector performance in VoIP- and telephony-like scenarios, while ReplayDF~\cite{muller2025replay} reports similar degradation under playback and re-recording. Related discussion in~\cite{delgado2026presentation} further points to a realism gap between raw generated audio and audio after practical presentation through communication channels or platform pipelines. These studies establish delivery as an important source of performance variation, but they typically examine communication, replay, or platform effects separately rather than as parts of one structured delivery process.

Recent frameworks have also improved reproducibility at the system level. AUDDT~\cite{zhu2025auddt} standardizes cross-dataset evaluation across public corpora, and Spoof-SUPERB~\cite{ali2026superb} provides a SUPERB-style comparison of self-supervised speech representations for audio deepfake detection. These efforts improve fair and reproducible comparison, but they primarily reorganize or evaluate existing resources rather than introduce a matched-parent benchmark centered on post-generation delivery.

ChainBench-ADD is designed to complement this landscape by modeling post-generation delivery as an ordered process under matched-parent control. To the best of our knowledge, existing representative resources have not jointly provided multiple post-generation delivery families, explicit ordered templates, and matched-parent local analysis within one unified benchmark.

\section{The ChainBench-ADD Benchmark}
\label{sec:benchmark}

\subsection{Overview and Terminology}
ChainBench-ADD organizes post-generation delivery into families, templates, variants, and realized chains. This hierarchy supports scenario-, route-, local-, and path-level analysis under matched-parent control.

Throughout this section, a delivery operator denotes an atomic or composite transform such as codec compression, re-encoding, packet loss, room reverberation, or additive noise. A template is an ordered operator blueprint within one delivery family, and a variant is one parameterized instantiation. A delivery chain is the realized ordered path applied to one parent sample. A delivery family is the top-level scenario category that groups related templates, including direct, platform-like, telephony, simulated replay, and hybrid delivery.

\subsection{Delivery Families and Chain Realization}
ChainBench-ADD defines five delivery families that reflect common post-generation routes. Direct is the identity control. Platform-like approximates online redistribution and recompression. Telephony models bandwidth-limited call transmission with telephony codecs, packet impairment, and compact end-to-end call simulation. Simulated replay models playback and recapture in physical spaces through room response and environmental noise. Hybrid combines communication and acoustic effects within a single path. These families are best understood as scenario categories rather than disjoint operator inventories: different families may share operators but combine them in different ordered patterns and parameter regimes.

These families are built from eight reusable delivery operators: resampling, band-limiting, codec compression, re-encoding, packet loss, additive noise, room impulse response (RIR), and call-path simulation. Many operators are shared across multiple families. For example, codec-like effects appear in both platform-like and telephony delivery, while RIR appears in both simulated replay and hybrid delivery. A template specifies an ordered operator blueprint within one delivery family, a variant instantiates that blueprint with concrete parameter choices, and a realized chain is the waveform-domain rendering of that ordered path for one parent sample. Table~\ref{tab:ChainBench-ADD-operators} summarizes the operator inventory and representative parameter ranges of each operator. Full parameter pools are deferred to the appendix.

\begin{table}[t]
\caption{Delivery operators in ChainBench-ADD.}
\label{tab:ChainBench-ADD-operators}
\vspace{-0.5em}
\centering
\footnotesize
\setlength{\tabcolsep}{2pt}
\renewcommand{\arraystretch}{1.30}
\begin{tabularx}{\linewidth}{@{}%
    >{\raggedright\arraybackslash}p{0.17\linewidth}%
    @{\hspace{2pt}}%
    >{\raggedright\arraybackslash}X%
    @{\hspace{2pt}}%
    >{\centering\arraybackslash}p{0.08\linewidth}%
    @{}}
\toprule
\textbf{Operator} & \textbf{Representative settings} & \textbf{Fam.} \\
\midrule
Resample 
& 16$\rightarrow$8, 16$\rightarrow$24, 16$\rightarrow$8$\rightarrow$16, 
16$\rightarrow$24$\rightarrow$16, 16$\rightarrow$32$\rightarrow$16 kHz 
& P/T/R/H \\

Band-limit 
& Narrowband, wideband 
& T/H \\

Codec 
& AAC (24/32/48 kbps), Opus (16/24/32 kbps), GSM, $\mu$-law PCM 
& P/T/H \\

Re-encode 
& Same/cross-codec AAC/Opus recompression (24/32 kbps) 
& P/R/H \\

Packet loss 
& 1/3/5/10\% loss; burst 2/3/5; repeat-fade, interpolation, or noise-fill concealment 
& T/H \\

Noise 
& White, pink, brown, babble, hiss, hum; 30/20/15/10 dB SNR 
& R \\

RIR 
& Small/medium/large rooms; RT60 0.2/0.4/0.6/0.8 s; source--microphone distance 0.5/1/2/3 m 
& R/H \\

Call-path 
& Joint channel filtering, codec, packet loss, jitter (0/8/16 ms), and AGC 
& T \\
\bottomrule
\end{tabularx}

\medskip
\raggedright
\footnotesize \textbf{Fam.} P = platform-like, T = telephony, R = simulated replay, H = hybrid.
\end{table}

All operators are implemented as waveform-domain transforms rather than symbolic tags. Codec and re-encode effects use actual encode--decode round trips, packet loss includes burstiness and concealment, RIR is applied with Pyroomacoustics~\cite{scheibler2018pyroomacoustics}, and call-path simulation jointly models channel filtering, codec compression, burst loss, jitter-buffer effects, and AGC. After rendering, all outputs are standardized to 16\,kHz mono PCM WAV.

This hierarchy supports analysis at multiple levels, including scenario-level robustness, generalization to unseen route patterns, operator-level sensitivity, and path-wise robustness under accumulated edits. The current release contains 33 templates in total (1 direct, 6 platform-like, 10 telephony, 7 simulated replay, and 9 hybrid), which realize 26 unique ordered operator sequences and are rendered into sample-specific realized chains. Figure~\ref{fig:chainbench-overview} summarizes this structure.

\begin{figure*}[t]
  \centering
  \captionsetup[subfigure]{justification=centering}
  \begin{subfigure}[t]{0.28\textwidth}
    \centering
    \includegraphics[width=\linewidth]{figures/chainbench_add_a_family_v5.pdf}
    \caption{Delivery-family composition.}
    \label{fig:chainbench-family}
  \end{subfigure}\hfill
  \begin{subfigure}[t]{0.34\textwidth}
    \centering
    \includegraphics[width=\linewidth]{figures/chainbench_add_b_family_operator_v5.pdf}
    \caption{Delivery-family/operator map.}
    \label{fig:chainbench-family-operator}
  \end{subfigure}\hfill
  \begin{subfigure}[t]{0.30\textwidth}
    \centering
    \includegraphics[width=\linewidth]{figures/chainbench_add_c_transition_v5.pdf}
    \caption{Operator transitions.}
    \label{fig:chainbench-transition}
  \end{subfigure}
  \caption{Composition and structure of ChainBench-ADD delivery chains. (a) Sample counts by delivery family, split into bona fide and spoof. (b) Template-level delivery-family/operator map, where each cell shows how many templates in a family contain the operator. (c) Adjacent-operator transition matrix induced by the template inventory.}
  \label{fig:chainbench-overview}
\end{figure*}

\subsection{Data Construction}
ChainBench-ADD is built with a staged pipeline that keeps transcript and speaker identity fixed while introducing spoofing and delivery effects in separate steps. This design follows the benchmark goal: to measure robustness to post-generation delivery artifacts rather than differences in content, speaker, or source formatting. We therefore first construct a clean parent pool, then generate spoof counterparts, and only afterward render delivery chains and export protocol-aware metadata. This ordering ensures that delivery is evaluated as a downstream factor rather than being conflated with generator choice, transcript content, or speaker identity.

We curate bona fide speech from Common Voice~24.0 for English~\cite{ardila2020commonvoicemassivelymultilingualspeech} and AISHELL-3~\cite{shi2020aishell} for Mandarin Chinese. Candidate utterances are screened at both the text and audio levels, removing unsuitable transcripts and low-quality recordings. Selection is speaker-centric: for each accepted speaker, a fixed number of utterances is retained after ranking candidates by properties such as usable duration, speech activity, and stable loudness. We then assign benchmark speaker IDs and create train/dev/test splits at the speaker level. The retained utterances are converted into clean masters by trimming silence, downmixing to mono, resampling to 16\,kHz, normalizing loudness, and saving as 16-bit PCM WAV, yielding 18,703 bona fide clean parents.

For each bona fide clean parent, we synthesize two spoof clean parents using two generators sampled from six contemporary systems: Qwen3-TTS~\cite{Qwen3-TTS}, CosyVoice3~\cite{du2025cosyvoice3}, Spark-TTS~\cite{wang2025sparktts}, F5-TTS~\cite{chen-etal-2024-f5tts}, IndexTTS2~\cite{zhou2025indextts2}, and VoxCPM~\cite{voxcpm2025}. The generator set is chosen to span multiple synthesis families, including multilingual instruction TTS, zero-shot cloning TTS, expressive controllable TTS, and architecture-diverse TTS. To preserve identity while avoiding trivial prompt leakage, the target transcript is paired with a different clean utterance from the same speaker as prompt/reference audio whenever possible. After validation, this stage yields 37,110 spoof clean parents and 55,813 clean parents in total.

Delivery chains are rendered only after the clean pool is fixed, keeping delivery orthogonal to the bona fide/spoof label. For each parent, we generate one direct control and sampled variants from the non-direct families by instantiating ordered waveform-domain operator chains. The pipeline also preserves paired constructions for controlled local comparisons, including order-swap protocols. Final packaging validates audio and exports metadata that links each sample to its parent, family, template, variant, and operator sequence. This explicit lineage is what later enables matched local interventions and path-wise robustness analysis under shared parent context.

\subsection{Benchmark Tasks and Protocols}
ChainBench-ADD is released as a metadata-driven benchmark for delivery-aware audio deepfake detection. From the shared metadata, we instantiate five tasks: in-chain detection, operator substitution, parameter perturbation, order swap, and delivery robustness.

The base task is in-chain detection, which performs bona fide-vs.-spoof classification separately within each delivery family. The three intervention tasks are constructed under matched parent, family, and label control. In operator substitution, paired samples have the same parent utterance, delivery family, ground-truth label, and chain length, but differ only in the operator identity at one or more positions. In parameter perturbation, paired samples differ at exactly one operator position while keeping the operator identity fixed and changing only its parameterization. In order swap, paired samples preserve the same operators and label but differ by one adjacent minimal swap in operator order. Finally, delivery robustness groups samples into lineage graphs defined by shared parent utterance, delivery family, and label, and evaluates whether correct decisions persist as additional observed delivery edits accumulate along the chain. Combined with an extra leave-one-template-out evaluation, these protocols test family-level robustness, transfer to unseen templates, sensitivity to controlled structural edits, and stability along realistic delivery paths.

\subsection{Dataset Statistics, Metadata, and Validation}
ChainBench-ADD contains 941,201 waveforms derived from 55,813 parents and 448 speakers across English and Mandarin Chinese, including 315,573 bona fide and 625,628 spoof samples. The speaker-disjoint train/dev/test split contains 663,363/136,027/141,811 samples from 314/67/67 speakers, respectively. At the delivery level, the release covers five delivery families, 33 templates, and 26 ordered operator sequences. At the protocol level, it further provides 182,409 operator-substitution groups, 53,234 parameter-perturbation groups, 108,245 order-swap groups, and 27,204 lineage/path groups.

The released metadata are lightweight and protocol-aware. In addition to identifiers, transcript, language, speaker, and provenance, they include split label, generator family and model, delivery family, template and variant identifiers, ordered operator signatures, and the structural annotations needed to reconstruct the benchmark tasks. Before export, each waveform is rechecked against target format and duration constraints, and the final package records split-integrity checks, duplicate-identifier and duplicate-path checks, and parent-coverage summaries across delivery families.

We further validate matched-parent consistency on the released 141,811-sample test split derived from 8,407 parents. Using a transformers ASR backend (Whisper-large-v3-turbo) and a transformers speaker-embedding backend (WavLM-base-sv), we compare each delivered sample against its clean parent through transcript error-rate drift, speaker cosine similarity, duration ratio, and absolute RMS drift. Full family-wise statistics are deferred to the appendix.
\section{Experiments}
\label{sec:experiments}

\subsection{Baselines and Training Setup}
We evaluate five representative baselines: AASIST~\cite{Jung2021AASIST}, AASIST-L~\cite{Jung2021AASIST}, XLS-R+SLS~\cite{sls}, SafeEar~\cite{li2024safeear}, and Nes2Net~\cite{Nes2Net}. Together, they cover graph-based anti-spoofing, lightweight anti-spoofing, self-supervised speech representations, content privacy-preserving detection, and recent efficient detector design. Our goal is not to exhaust the detector literature, but to test whether the benchmark reveals delivery-conditioned patterns that are consistent across heterogeneous model families.

We use the official implementations of all baselines and modify only data loading and score export for ChainBench-ADD. AASIST and AASIST-L are trained for 20 epochs with batch size 24 and learning rate $10^{-3}$. XLS-R+SLS, SafeEar, and Nes2Net are trained for 15 epochs using their model-specific batch sizes and learning rate $10^{-6}$. Weight decay is set to $10^{-4}$ for all models.


\subsection{Evaluation Metrics and Settings}
We follow the five benchmark protocols defined in Sec.~\ref{sec:benchmark} and focus here on evaluation setup and metrics.

\noindent\textbf{In-chain detection}
We report family-wise equal error rate (EER, \%) for bona fide-vs.-spoof classification within each delivery family, together with the macro average across families. To assess cross-template generalization, we further adopt a leave-one-template-out protocol for each non-direct family, excluding one template during training and using it only for testing; the reported EER is averaged over held-out templates within that family.

\noindent\textbf{Intervention robustness}
For operator substitution, parameter perturbation, and order swap, we evaluate matched test pairs using the same train/dev pools as in the corresponding benchmark task. We report two primary metrics. The first is pair consistency, i.e., the proportion of pairs for which the model gives the same binary decision to both samples. The second is both-correct rate, i.e., the proportion of pairs for which both samples are classified correctly. Together, these metrics show whether stability under a controlled intervention reflects correct discrimination rather than repeated error. For all decision-based robustness metrics, we use a single EER-defined reference threshold computed once on the pooled evaluation split of each task and keep it fixed for all matched-pair subsets and lineage graphs, without any family- or subgroup-specific retuning. As a secondary score-level diagnostic, we also report the mean normalized score difference (MNSD) between the two samples in each pair. These decision-based metrics should therefore be interpreted as standardized diagnostics at a common operating point rather than as deployment-calibrated results. Higher values are better for pair consistency and both-correct rate, while lower values are better for MNSD. Formal definition details are provided in the appendix.

\noindent\textbf{Path-wise delivery robustness}
We evaluate delivery robustness on lineage graphs. For each lineage, the reference node is chosen as the shortest observed operator-signature realization, and graph distance is measured by shortest path over observed atomic delivery edits. We report the censored first-failure distance (C-FFD), which measures how far a detector remains correct before the first error appears and assigns a maximal value when no failure is observed in the evaluation set. We also compute the fraction of lineages that remain correct within $k$ accumulated delivery edits, denoted by $R_k$, and summarize this robustness curve by its area under the curve (AURC-chain). Larger values indicate more persistent correctness under accumulated delivery edits. Their Formal definitions are also provided in the appendix.


\subsection{Results and Analysis}

\noindent\textbf{Delivery family strongly affects detection difficulty.}
Table~\ref{tab:exp_inchain} reports in-chain detection EER by delivery family. The difficulty of audio deepfake detection varies systematically with the delivery family. Averaged across baselines, direct is the easiest condition (14.80 EER), followed by platform-like (19.12) and telephony (24.81), whereas simulated replay (31.77) and hybrid (31.91) are consistently more difficult. Replay-like recapture and mixed delivery effects introduce more degradation than mild redistribution. Nes2Net and XLS-R+SLS obtain the lowest macro EERs, but all baselines follow the same broad pattern across families. This reveals a stable delivery-conditioned difficulty structure that pooled evaluation would obscure.

\begin{table}[t]
\centering
\footnotesize
\setlength{\tabcolsep}{4pt}
\renewcommand{\arraystretch}{1.02}
\caption{In-chain detection EER (\%) by delivery family. Lower is better.}
\label{tab:exp_inchain}
\vspace{-0.5em}
\begin{tabular}{@{}lcccccc@{}}
\toprule
\textbf{Model} & \textbf{Direct} & \textbf{Platform} & \textbf{Telephony} & \textbf{SimReplay} & \textbf{Hybrid} & \textbf{Macro} \\
\midrule
AASIST       & 13.00 & 24.06 & 27.29 & 39.23 & 42.26 & 29.17 \\
AASIST-L     & 16.89 & 24.49 & 29.30 & 41.57 & 38.79 & 30.21 \\
XLS-R+SLS    & \textbf{12.92} & 11.04 & 16.01 & 20.90 & \textbf{21.11} & 16.40 \\
SafeEar      & 18.14 & 27.43 & 35.65 & 37.93 & 35.94 & 31.02 \\
Nes2Net      & 13.03 & \textbf{8.57} & \textbf{15.82} & \textbf{19.21} & 21.44 & \textbf{15.61} \\
\bottomrule
\end{tabular}
\end{table}

\noindent\textbf{The delivery-family ordering persists under unseen templates.}
Table~\ref{tab:exp_holdout} reports leave-one-template-out in-chain detection. The overall pattern remains close to the standard in-chain results: platform-like is still the easiest family, telephony remains intermediate, and simulated replay/hybrid remain the most difficult. The family ordering therefore appears to reflect effects beyond template memorization and transfers to unseen delivery realizations.

\begin{table}[t]
\centering
\footnotesize
\setlength{\tabcolsep}{4pt}
\renewcommand{\arraystretch}{1.02}
\caption{Leave-one-template-out in-chain detection EER (\%). Lower is better.}
\label{tab:exp_holdout}
\vspace{-0.5em}
\begin{tabular*}{0.96\linewidth}{@{\extracolsep{\fill}}lccccc@{}}
\toprule
\textbf{Model} & \textbf{Platform} & \textbf{Telephony} & \textbf{SimReplay} & \textbf{Hybrid} & \textbf{Macro} \\
\midrule
AASIST       & 23.52 & 28.52 & 41.66 & 41.55 & 33.81 \\
AASIST-L     & 26.53 & 29.04 & 40.68 & 42.27 & 34.63 \\
XLS-R+SLS    & 13.79 & 16.14 & 21.41 & 24.03 & 18.84 \\
SafeEar      & 28.96 & 33.92 & 39.49 & 38.21 & 35.15 \\
Nes2Net      & \textbf{11.33} & \textbf{16.08} & \textbf{21.32} & \textbf{24.09} & \textbf{18.20} \\
\bottomrule
\end{tabular*}
\end{table}

\noindent\textbf{Matched local interventions reveal different failure modes.}
Table~\ref{tab:exp_pair} shifts the focus from family-level variation to controlled local edits. The included three intervention tasks stress detectors in different ways. Operator substitution is the most disruptive intervention overall: averaged across baselines, it yields the lowest both-correct rate and the largest MNSD, indicating that changing operator identity perturbs detector behavior more than changing a parameter within the same operator. Parameter perturbation is the mildest intervention on average, suggesting that many detectors retain partial tolerance to within-operator variation even when they are less robust to larger structural changes in the chain. Order swap exposes a different weakness. It has the highest average pair consistency of the three tasks, yet its average both-correct rate remains well below its consistency level, especially for weaker baselines. The gap between these two quantities is largest for AASIST, AASIST-L, and SafeEar, which suggests that local reordering often preserves a model's decision without preserving its correctness. In other words, these models can be wrong on both elements of a matched pair. By contrast, XLS-R+SLS and Nes2Net show much smaller consistency--correctness gaps across all three interventions, suggesting that their invariance is more often aligned with correct discrimination rather than merely repeated error. Taken together, the matched-paired tasks help distinguish between sensitivity to operator identity, tolerance to moderate parameter variation, and vulnerability to local sequence changes. This failure-mode information is difficult to obtain from the pooled detection metrics alone.

\begin{table}[t]
\centering
\scriptsize
\setlength{\tabcolsep}{3pt}
\renewcommand{\arraystretch}{1.03}
\caption{Matched-pair robustness. Cons.=pair consistency, Both=both-correct rate, and MNSD=mean normalized score difference. Higher is better for Cons./Both; lower is better for MNSD.}
\label{tab:exp_pair}
\vspace{-0.5em}
\begin{tabular}{@{}lccccccccc@{}}
\toprule
\textbf{} & \multicolumn{3}{c}{\textbf{Operator substitution}} & \multicolumn{3}{c}{\textbf{Parameter perturbation}} & \multicolumn{3}{c}{\textbf{Order swap}} \\
\cmidrule(lr){2-4} \cmidrule(lr){5-7} \cmidrule(lr){8-10}
\textbf{Model} & \textbf{Cons.}$\bm{\uparrow}$ & \textbf{Both}$\bm{\uparrow}$ & \textbf{MNSD}$\bm{\downarrow}$ & \textbf{Cons.}$\bm{\uparrow}$ & \textbf{Both}$\bm{\uparrow}$ & \textbf{MNSD}$\bm{\downarrow}$ & \textbf{Cons.}$\bm{\uparrow}$ & \textbf{Both}$\bm{\uparrow}$ & \textbf{MNSD}$\bm{\downarrow}$ \\
\midrule
AASIST     & 74.80 & 54.79 & 52.67 & 78.16 & 62.41 & 43.49 & 80.98 & 58.74 & 41.20 \\
AASIST-L   & 73.75 & 52.97 & 53.10 & 77.22 & 60.88 & 34.66 & 79.19 & 55.56 & 35.16 \\
XLS-R+SLS  & \textbf{87.74} & 81.07 & 20.47 & 90.36 & 85.99 & 16.48 & \textbf{90.48} & \textbf{82.62} & 16.47 \\
SafeEar    & 65.25 & 47.60 & 52.90 & 67.40 & 51.29 & 45.69 & 72.56 & 51.23 & 39.85 \\
Nes2Net    & 87.57 & \textbf{81.86} & \textbf{15.47} & \textbf{90.45} & \textbf{86.92} & \textbf{11.70} & 89.32 & 81.83 & \textbf{13.72} \\
\bottomrule
\end{tabular}
\end{table}

\noindent\textbf{Robustness along delivery paths adds information beyond isolated edits.}
Table~\ref{tab:exp_delivery} evaluates whether correctness persists as delivery edits accumulate along realistic lineage paths. Unlike the matched-pair tasks, which isolate a single local change, this protocol tests stability under multi-step variation. The lineage metrics further separate the two strongest models. Although XLS-R+SLS and Nes2Net are similarly strong under isolated interventions, Nes2Net is consistently better in C-FFD, AURC-chain, and all reported $R_k$ values, suggesting more persistent correctness as delivery edits accumulate. The radius-wise results also show that most newly induced errors appear early: for every model, the largest drop is from $R_0$ to $R_1$, after which the curves flatten substantially, especially for the stronger baselines. Path-wise evaluation therefore shows both when failures first emerge and how quickly they propagate, complementing the view provided by isolated matched-pair tests.

\begin{table}[t]
\centering
\footnotesize
\setlength{\tabcolsep}{4pt}
\renewcommand{\arraystretch}{1.02}
\caption{Lineage-level delivery robustness. C-FFD denotes censored first-failure distance. Higher is better for all metrics.}
\label{tab:exp_delivery}
\vspace{-0.5em}
\begin{tabular*}{0.96\linewidth}{@{\extracolsep{\fill}}lcccccc@{}}
\toprule
\textbf{Model} & \textbf{C-FFD} & \textbf{AURC-chain} & \textbf{$\bm{R_0}$} & \textbf{$\bm{R_1}$} & \textbf{$\bm{R_2}$} & \textbf{$\bm{R_3}$} \\
\midrule
AASIST     & 2.79 & 69.73 & 73.82 & 68.60 & 68.25 & 68.24 \\
AASIST-L   & 2.61 & 65.20 & 69.21 & 64.09 & 63.76 & 63.74 \\
XLS-R+SLS  & 3.47 & 86.76 & 88.39 & 86.31 & 86.18 & 86.18 \\
SafeEar    & 2.46 & 61.45 & 66.97 & 59.93 & 59.45 & 59.44 \\
Nes2Net    & \textbf{3.51} & \textbf{87.86} & \textbf{89.52} & \textbf{87.40} & \textbf{87.27} & \textbf{87.26} \\
\bottomrule
\end{tabular*}
\end{table}

% \noindent\textbf{Summary.}
% The experiments show that ChainBench-ADD captures complementary aspects of delivery-aware audio deepfake detection. Delivery family induces a stable difficulty hierarchy, with direct/platform-like conditions generally easier and simulated replay/hybrid conditions harder. This ordering largely survives leave-one-template-out evaluation, suggesting family-level effects beyond template memorization. The matched intervention tasks reveal distinct local failure modes: operator substitution is typically the most disruptive edit, parameter perturbation is milder, and order swap often separates decision consistency from actual correctness. Lineage-based evaluation adds the cumulative view, showing that many failures emerge within the first one or two edits and that models with similar local robustness can diverge under multi-step delivery variation.

\section{Conclusion}
\label{sec:conclusion}

Current audio deepfake benchmarks mainly emphasize the generation side, while post-generation delivery is less systematically represented. We introduce ChainBench-ADD, a delivery-aware dataset and benchmark that organizes post-generation delivery through structured families, reusable operators, ordered templates, and realized chains under matched-parent control. Experiments with representative baselines show that delivery remains a major challenge: replay-like and hybrid routes are consistently harder than direct or lightly processed conditions, and protocol-aware evaluation reveals robustness patterns that pooled testing alone can miss. We hope ChainBench-ADD can serve as a standardized and reproducible benchmark for delivery-aware evaluation, and support future research on robust audio deepfake detection under realistic deployment conditions.

% \clearpage

%%
%% The acknowledgments section is defined using the "acks" environment
%% (and NOT an unnumbered section). This ensures the proper
%% identification of the section in the article metadata, and the
%% consistent spelling of the heading.
% \begin{acks}
% To Robert, for the bagels and explaining CMYK and color spaces.
% \end{acks}

%%
%% The next two lines define the bibliography style to be used, and
%% the bibliography file.
\bibliographystyle{ACM-Reference-Format}
\bibliography{main}


%%
%% If your work has an appendix, this is the place to put it.
% \clearpage
% \appendix

\section{Delivery-chain renderer specification}
\label{app:renderer_spec}

We summarize the released renderer at four levels: family-level sampling policy, template inventory, fixed parameter pools, and waveform-level operator realization. Unless noted otherwise, the appendix tables are authoritative for the exact configuration used to build the benchmark package in this paper. The operator summary in the main text is intentionally scenario-level; when it is less specific than the appendix, the appendix should be read as the exact release description.

\subsection{Family-level rendering logic}

For each parent utterance, the renderer first selects a delivery family, then samples a concrete template from that family, and finally instantiates the template-level operators from fixed parameter pools. The direct family contributes a single identity-control child, whereas each non-direct family contributes four rendered children. Within a family, template identifiers are deterministically shuffled per parent and sampled without replacement. For paired templates, one canonical member of the pair is instantiated first and its realized operators are then reordered to form the paired counterpart, so that the two children differ only in local order while sharing the same sampled parameters.

\begin{table*}[t]
\caption{Family-level rendering policy in the released renderer.}
\label{tab:appendix_family_policy}
\centering
\small
\setlength{\tabcolsep}{5pt}
\renewcommand{\arraystretch}{1.10}
\begin{tabular}{@{}p{0.14\textwidth}p{0.10\textwidth}p{0.14\textwidth}p{0.42\textwidth}@{}}
\toprule
\textbf{Family} & \textbf{Templates} & \textbf{Children / parent} & \textbf{Sampling rule} \\
\midrule
Direct & 1 & 1 & Identity control only. \\
Platform-like & 6 & 4 & Sample without replacement; no paired-template forcing. \\
Telephony & 10 & 4 & Sample without replacement; no paired-template forcing. \\
Simulated replay & 7 & 4 & One paired group may be emitted first; otherwise sample without replacement. \\
Hybrid & 9 & 4 & One paired group may be emitted first; otherwise sample without replacement. \\
\bottomrule
\end{tabular}
\end{table*}

A further detail concerns re-encoding. In same mode, the re-encode operator is forced to use the most recent codec already realized earlier in the chain. In cross mode, it is forced to differ from that codec and is sampled from \{AAC, Opus\}. If no earlier codec exists, the renderer falls back to a family-specific default: AAC for simulated replay and Opus for the relevant hybrid templates. In addition, whenever the realized codec is GSM, the encoding sample rate is fixed to 8~kHz and decoding also proceeds from 8~kHz.

\subsection{Template inventory and parameter pools}

The configuration used here contains 33 templates, which realize 26 distinct ordered operator sequences because several templates intentionally keep the same operator multiset while changing only local order. Table~\ref{tab:appendix_templates} gives the exact family--template--operator inventory used in the release. In this inventory, call\_path appears only in the three telephony session templates (session\_nb\_mulaw, session\_nb\_gsm, and session\_wb\_opus), whereas reencode appears in platform-like, simulated replay, and hybrid, but not telephony.

\begin{table*}[t]
\caption{Template inventory as ordered operator blueprints before parameter instantiation.}
\label{tab:appendix_templates}
\centering
\footnotesize
\setlength{\tabcolsep}{2pt}
\renewcommand{\arraystretch}{1.04}
\begin{tabular}{@{}p{0.11\textwidth}p{0.18\textwidth}p{0.27\textwidth}@{}}
\toprule
\textbf{Family} & \textbf{Template ID} & \textbf{Ordered operator sequence} \\
\midrule
Direct & direct\_clean & Identity chain with no operator. \\
\midrule
\multirow[t]{6}{*}{Platform-like} & aac\_single & codec \\
& opus\_single & codec \\
& aac\_reencode & codec $\rightarrow$ reencode \\
& opus\_reencode & codec $\rightarrow$ reencode \\
& aac\_resample\_reencode & codec $\rightarrow$ resample $\rightarrow$ reencode \\
& resample\_opus & resample $\rightarrow$ codec \\
\midrule
\multirow[t]{10}{*}{Telephony} & nb\_mulaw & bandlimit $\rightarrow$ codec \\
& nb\_gsm & bandlimit $\rightarrow$ codec \\
& wb\_opus & bandlimit $\rightarrow$ codec \\
& nb\_mulaw\_plr & bandlimit $\rightarrow$ codec $\rightarrow$ packet\_loss \\
& nb\_resample\_mulaw\_plr & resample $\rightarrow$ bandlimit $\rightarrow$ codec $\rightarrow$ packet\_loss \\
& wb\_resample\_opus\_plr & resample $\rightarrow$ bandlimit $\rightarrow$ codec $\rightarrow$ packet\_loss \\
& wb\_opus\_resample\_return & bandlimit $\rightarrow$ codec $\rightarrow$ resample \\
& session\_nb\_mulaw & call\_path \\
& session\_nb\_gsm & call\_path \\
& session\_wb\_opus & call\_path \\
\midrule
\multirow[t]{7}{*}{Simulated replay} & rir\_only & rir \\
& rir\_noise & rir $\rightarrow$ noise \\
& noise\_rir & noise $\rightarrow$ rir \\
& rir\_reencode & rir $\rightarrow$ reencode \\
& reencode\_rir & reencode $\rightarrow$ rir \\
& rir\_noise\_resample & rir $\rightarrow$ noise $\rightarrow$ resample \\
& resample\_rir\_reencode & resample $\rightarrow$ rir $\rightarrow$ reencode \\
\midrule
\multirow[t]{9}{*}{Hybrid} & opus\_plr\_rir & codec $\rightarrow$ packet\_loss $\rightarrow$ rir \\
& rir\_aac & rir $\rightarrow$ codec \\
& aac\_rir & codec $\rightarrow$ rir \\
& bandlimit\_codec\_rir & bandlimit $\rightarrow$ codec $\rightarrow$ rir \\
& rir\_bandlimit\_codec & rir $\rightarrow$ bandlimit $\rightarrow$ codec \\
& rir\_reencode\_plr & rir $\rightarrow$ reencode $\rightarrow$ packet\_loss \\
& reencode\_rir\_plr & reencode $\rightarrow$ rir $\rightarrow$ packet\_loss \\
& resample\_codec\_rir & resample $\rightarrow$ codec $\rightarrow$ rir \\
& bandlimit\_resample\_codec & bandlimit $\rightarrow$ resample $\rightarrow$ codec \\
\bottomrule
\end{tabular}
\end{table*}

Table~\ref{tab:appendix_parameter_pools} summarizes only the pool ranges and export contract. The exact per-family realizations are already determined by the template inventory and the operator descriptions below.

\begin{table*}[t]
\caption{Fixed parameter pools and export constraints in the released configuration.}
\label{tab:appendix_parameter_pools}
\centering
\small
\setlength{\tabcolsep}{5pt}
\renewcommand{\arraystretch}{1.10}
\begin{tabular}{@{}p{0.18\textwidth}p{0.50\textwidth}p{0.16\textwidth}@{}}
\toprule
\textbf{Group} & \textbf{Released range or set} & \textbf{Used by} \\
\midrule
Codec bitrates & AAC \{24, 32, 48\}; Opus \{16, 24, 32\}; telephony Opus \{16, 24\}; re-encode \{24, 32\} kbps & codec families \\
Resample modes & Family-specific one-way and round-trip 8/16/24/32 kHz conversions & resample \\
Re-encode codec set & AAC or Opus, with same/cross selection & re-encode \\
Packet-loss settings & Loss \{1, 3, 5, 10\}\%; burst \{2, 3, 5\} frames; three concealment modes & packet loss, call path \\
Call-path settings & Jitter \{0, 8, 16\} ms; AGC \{mild, telephony\} & call path \\
Noise settings & Six synthetic noise types; SNR \{30, 20, 15, 10\} dB & noise \\
RIR settings & RT60 \{0.2, 0.4, 0.6, 0.8\} s; distance \{0.5, 1.0, 2.0, 3.0\} m; three room presets & RIR \\
Export contract & Mono 16 kHz PCM WAV; duration checked in [1, 30] s before export & all outputs \\
\bottomrule
\end{tabular}
\end{table*}

\subsection{Operator realization}

All operators are realized as waveform-domain transforms rather than symbolic metadata tags, and every realization is recorded in the manifest trace.

\noindent\textbf{Resampling} The resampling operator uses ffmpeg with the SoX resampler at precision 28. One-way modes include 16k$\rightarrow$8k, 16k$\rightarrow$24k, 16k$\rightarrow$32k, 8k$\rightarrow$16k, 24k$\rightarrow$16k, and 32k$\rightarrow$16k. Round-trip modes are implemented as two successive resampling steps with an intermediate temporary waveform.

\noindent\textbf{Band-limiting} The band-limit operator is implemented with ffmpeg filtering. The narrowband profile applies a second-order high-pass filter at 250~Hz and a second-order low-pass filter at 3400~Hz, followed by companding with 0.01~s attack and 0.15~s decay. The wideband profile applies a second-order high-pass filter at 50~Hz and a second-order low-pass filter at 7000~Hz. To avoid terminology drift, we use \emph{narrowband (250--3400~Hz)} and \emph{wideband (50--7000~Hz)} consistently throughout the appendix and treat these implementation-level cutoffs as authoritative for the released benchmark.

\noindent\textbf{Codec and re-encoding} The codec and re-encode operators both perform true encode--decode round trips through ffmpeg. The supported codecs are AAC, Opus, GSM, $\mu$-law PCM, and A-law PCM. After encoding to the corresponding intermediate container, the audio is decoded back to mono PCM WAV. The only difference between these two operators is how the target codec is chosen: codec follows the template specification directly, while reencode follows the context-aware same/cross rule described above.

\noindent\textbf{Packet loss} The packet-loss operator works on 20~ms waveform frames. Let $p$ denote the target loss rate after clamping to at most 0.95 and let $b$ denote the requested average burst length in frames. The underlying two-state good/bad process uses
\begin{equation}
P(\mathrm{bad}\rightarrow\mathrm{good}) = \min(1, 1/b)
\label{eq:plr_bad_to_good}
\end{equation}
and
\begin{equation}
P(\mathrm{good}\rightarrow\mathrm{bad}) = \min\!\left(1, \frac{p\,P(\mathrm{bad}\rightarrow\mathrm{good})}{1-p}\right).
\label{eq:plr_good_to_bad}
\end{equation}
When a frame enters the bad state, one of three concealment strategies is applied: repeat-fade, interpolation, or noise fill. The processed frames are then concatenated back into a waveform.

\noindent\textbf{Room impulse responses} The RIR operator first attempts room simulation with Pyroomacoustics. For a requested room preset and RT60, it computes an absorption coefficient by inverse Sabine, samples valid source and microphone positions, generates a room impulse response, truncates the late tail, and convolves it with the waveform. If Pyroomacoustics is unavailable or fails, the renderer falls back to a synthetic RIR generator based on a direct path, exponentially decaying random reflections, and a low-level noise floor. The reverberant output is peak-normalized after convolution.

\noindent\textbf{Additive noise} The noise operator adds synthetic noise at a target SNR. White noise uses i.i.d. Gaussian samples; pink noise uses $1/\sqrt{f}$ spectral shaping; brown noise is obtained by cumulative summation; hiss emphasizes high frequencies; hum combines 50, 100, and 150~Hz tonal components with low-level white noise; and babble is synthesized by summing six randomly shifted and scaled noise streams. For babble, hum, and hiss, a non-stationary amplitude envelope is applied before SNR calibration.

\noindent\textbf{Call paths} The call-path operator is a composite transform with a fixed internal pipeline: profile-specific band filtering and resampling to the target sample rate, codec round trip, burst packet loss with the same implementation as above, frame-wise jitter-buffer simulation, and AGC or limiting. The wideband profile uses the same 50--7000~Hz passband as the standalone band-limit operator. The narrowband profile uses the same 250~Hz and 3400~Hz cutoffs but a different compand curve. Jitter is implemented by displacing each frame by a uniformly sampled offset in $[-J, J]$, overlap-averaging the displaced frames, and trimming or padding the result back to the original length.

\section{Released metadata \& protocol reconstruction}
\label{app:release_metadata}

We now turn from waveform generation to how the packaged benchmark is encoded. The release contains one global metadata CSV and split-specific metadata CSV files. Each row describes one waveform that survives Stage-5 validation and export. The same packaging step also writes a dataset summary with speaker-disjointness checks, duplicate-ID and duplicate-path checks, dropped-row logs, and parent-coverage summaries. Together, these files form the executable interface for reconstructing the benchmark protocols.

Table~\ref{tab:appendix_manifest_schema} summarizes the protocol-critical field groups. In addition to standard identifiers and provenance, the metadata records the ordered operator sequence, the order-invariant operator multiset, realized delivery parameters such as codec, bitrate, packet loss, bandwidth mode, SNR, RT60, room dimensions, source--microphone distance, and the instantiation seed.

\begin{table*}[t]
\caption{Protocol-critical released metadata groups.}
\label{tab:appendix_manifest_schema}
\centering
\small
\setlength{\tabcolsep}{5pt}
\renewcommand{\arraystretch}{1.08}
\begin{tabular}{@{}p{0.20\textwidth}p{0.56\textwidth}@{}}
\toprule
\textbf{Field group} & \textbf{Why it matters} \\
\midrule
Identity and paths & Links each delivered child to its clean parent and exported waveform. \\
Labels and provenance & Supports standard train/dev/test use, filtering, and speaker-disjoint checks. \\
Generator metadata & Enables provenance-aware analysis by generator family or system. \\
Delivery-chain metadata & Encodes the realized family, template, operator order, and realized parameters. \\
Matched-pair annotations & Reconstructs operator substitution, parameter perturbation, and order swap tasks. \\
Lineage/path annotations & Reconstructs delivery lineages and graph-distance based analyses. \\
Audio-format checks & Verifies that released files satisfy the final export contract. \\
\bottomrule
\end{tabular}
\end{table*}

The structural groups are derived deterministically from the metadata rather than manually enumerated. Operator-substitution groups fix the parent utterance, binary label, delivery family, and local chain context and require exactly one operator identity change. Parameter-perturbation groups fix operator identity and order and vary exactly one parameter axis. Order-swap groups fix the realized operator multiset and differ by one adjacent transposition. Path/lineage groups fix the parent utterance, binary label, and delivery family, then derive a shortest-signature reference node and graph-distance annotations from the observed realized chains. The exact string form of the group identifiers is implementation-defined, but the grouping semantics are fixed by these rules.


\section{Matched-parent preservation validation}
\label{app:parent_child_drift}

We validate the benchmark's matched-parent assumption: delivery rendering should alter surface acoustics without materially changing utterance content or speaker identity, so that the clean parent remains a valid reference for controlled comparison. The analysis runs from the packaged Stage-5 metadata and, by default, scans the exported test split. For each row it resolves the child waveform, its paired clean parent, and the canonical transcript, then writes row-level outputs keyed by sample\_id and parent\_id.

Each child--parent pair yields three preservation signals and two bookkeeping outputs. The script computes parent and child ASR hypotheses and measures transcript preservation against the canonical transcript using WER and CER, reporting $\Delta$WER and $\Delta$CER as child-minus-parent drift. It also embeds the parent and child waveforms with a fixed speaker backend and reports cosine similarity, while the waveform itself provides lightweight acoustic drift measures such as duration ratio, RMS change, and peak amplitude. In addition, the export records status and error fields and emits both row-level and aggregated summaries.

The text normalization and error-rate computation follow the released implementation. For English, the reference and hypothesis are lowercased, punctuation is stripped, and whitespace is collapsed before tokenization. For Chinese, the normalization keeps CJK and alphanumeric characters, and WER falls back to character-level tokenization when no whitespace segmentation is present. CER is always computed on the normalized character sequence. The row-level export retains both raw and normalized hypotheses so the preservation statistics can be recomputed exactly.

The validation backends are fixed. ASR uses the default transformers\_asr backend with a Hugging Face automatic-speech-recognition pipeline instantiated from openai/whisper-large-v3-turbo, chunk length 15\,s, batch size 16, generation task transcribe, and half-precision model loading. Speaker preservation uses the default transformers\_speaker backend, which loads a wavlm-base-sv x-vector checkpoint through AutoFeatureExtractor and AutoModelForAudioXVector, embeds parent and child audio in batches of 16, $\ell_2$-normalizes the embeddings, and computes cosine similarity. These backends are held fixed across all families and languages.

Overall, the mean child-minus-parent drift is 0.061 in WER and 0.043 in CER, with mean speaker cosine 0.941, mean duration ratio 1.056, and mean absolute RMS drift of 2.334\,dB. Platform-like remains closest to the parent, telephony introduces moderate degradation, and simulated replay and hybrid induce the largest shifts. This pattern supports the benchmark design: delivery transforms are strong enough to create a meaningful robustness challenge while still preserving the matched-parent structure required for controlled evaluation.

Table~\ref{tab:appendix_parent_child_drift} reports only family means because this is the clearest compact summary for the paper. The released preservation outputs additionally provide row-level records, family-language aggregates, and explicit status counts. We therefore aggregate bona fide and spoof rows together in the appendix: the preservation question concerns whether the delivery renderer preserves the parent-reference relationship, which is a property of the transformation pipeline rather than of the detection label. The row-level CSV retains the original label field, so median, percentile, or label-conditioned summaries can be reproduced directly from the released outputs without rerendering the dataset.

\begin{table}[t]
\centering
\footnotesize
\setlength{\tabcolsep}{4pt}
\renewcommand{\arraystretch}{1.02}
\caption{Matched-parent preservation statistics by delivery family. Entries are family means over successful child--parent comparisons. $\Delta$WER and $\Delta$CER denote child-minus-parent transcript error drift against the reference transcript; lower is better. Speaker cosine is higher-is-better. Duration ratio and $|\bm{\Delta\mathrm{RMS}}|$ summarize acoustic drift.}
\label{tab:appendix_parent_child_drift}
\begin{tabular*}{0.96\linewidth}{@{\extracolsep{\fill}}lccccc@{}}
\toprule
\textbf{Family} & \textbf{$\bm{\Delta}$WER} & \textbf{$\bm{\Delta}$CER} & \textbf{Speaker cos.} & \textbf{Dur. ratio} & \textbf{$\bm{|\Delta\mathrm{RMS}|}$} \\
\midrule
Direct    & 0.000 & 0.000 & 1.000 & 1.000 & 0.000 \\
Platform  & 0.012 & 0.011 & 0.986 & 1.005 & 0.245 \\
Telephony & 0.059 & 0.046 & 0.943 & 1.001 & 1.620 \\
SimReplay & 0.094 & 0.063 & 0.909 & 1.120 & 4.247 \\
Hybrid    & 0.095 & 0.065 & 0.910 & 1.114 & 3.842 \\
\bottomrule
\end{tabular*}
\end{table}

\section{Benchmark-specific robustness metrics}
\label{app:robust_metrics}

Here we define only the new robustness metrics introduced in this work. Standard detection metrics such as EER, AUC, accuracy, and $F_1$ are omitted for brevity. Throughout the appendix we use the full name \emph{mean normalized score drift} (MNSD) consistently; the abbreviated table header \emph{NSD} in the main experiments section refers to the same quantity.

\begin{table*}[t]
\caption{Compact summary of the benchmark-specific metrics.}
\label{tab:appendix_new_metrics_summary}
\centering
\small
\setlength{\tabcolsep}{5pt}
\renewcommand{\arraystretch}{1.10}
\begin{tabular}{@{}p{0.13\textwidth}p{0.11\textwidth}p{0.19\textwidth}p{0.10\textwidth}@{}}
\toprule
\textbf{Setting} & \textbf{Metric} & \textbf{What it captures} & \textbf{Better} \\
\midrule
Matched pairs & PCR & Decision stability & Higher \\
Matched pairs & PJA & Stable correctness & Higher \\
Matched pairs & MNSD & Score sensitivity & Lower \\
Matched pairs & SMR (aux.) & Pairwise error risk & Lower \\
Lineages & C-FFD & First-failure distance & Higher \\
Lineages & $R_k$ & Correctness up to distance $k$ & Higher \\
Lineages & AURC-chain & Overall lineage robustness & Higher \\
\bottomrule
\end{tabular}
\end{table*}

\subsection{Notation and thresholding}

Let each evaluation sample have score $s_i \in \mathbb{R}$ and binary label $y_i \in \{0,1\}$, where $y_i=1$ denotes bona fide speech and $y_i=0$ denotes spoofed speech. Larger scores indicate stronger bona fide evidence. For each benchmark task, let $\mathcal{E}$ denote the full scored evaluation split used for method comparison (in our experiments, the pooled test split for that task), with score set $S_{\mathcal{E}}$ and labels $Y_{\mathcal{E}}$. We define a single \emph{reference operating point} once per task by
\begin{equation}
\tau_{\mathrm{ref}} =
\begin{cases}
\tau_{\mathrm{EER}}(S_{\mathcal{E}}, Y_{\mathcal{E}}), & \text{if both classes are present in } \mathcal{E},\\
0.5, & \text{otherwise}.
\end{cases}
\label{eq:pair_threshold}
\end{equation}
The matched-pair and lineage metrics then convert scores to hard decisions using this same fixed threshold:
\begin{equation}
\hat{y}_i = \mathbb{I}[s_i \ge \tau_{\mathrm{ref}}].
\label{eq:pair_decision}
\end{equation}

Two implementation choices matter. First, $\tau_{\mathrm{ref}}$ is estimated \emph{once} from the full evaluation split $\mathcal{E}$ and then reused for all matched-pair sets, delivery families, and lineage analyses within that task; it is not recomputed for individual subsets. Second, the hard-decision robustness metrics are intended as \emph{standardized analysis summaries at a common evaluation-defined operating point}, not as a claim of deployment-ready threshold calibration. This avoids subgroup-specific retuning and anchors all robustness comparisons to the same decision rule.

Accordingly, the threshold-free metrics (notably EER/AUC in the main experiments and MNSD below) should be read as the primary calibration-independent summaries, while PCR, PJA, C-FFD, $R_k$, and AURC-chain are decision-level diagnostics at the common reference operating point $\tau_{\mathrm{ref}}$. The use of $\tau_{\mathrm{EER}}$ therefore standardizes the binary robustness analyses rather than optimizing each protocol independently.

\subsection{Matched-pair metrics}

For operator substitution, parameter perturbation, and order swap, evaluation is carried out on the set of valid matched pairs
\begin{equation}
\mathcal{P} = \{(i_m,j_m)\}_{m=1}^{M},
\label{eq:pair_set}
\end{equation}
where the two members of each pair share the same ground-truth label and differ only in the targeted local intervention.

\noindent\textbf{Pair consistency rate} The pair consistency rate is defined as
\begin{equation}
\mathrm{PCR} =
\frac{1}{M}
\sum_{m=1}^{M}
\mathbb{I}[\hat{y}_{i_m}=\hat{y}_{j_m}],
\label{eq:pcr}
\end{equation}
and measures whether the detector preserves its binary decision under the controlled local edit, regardless of whether that decision is correct.

\noindent\textbf{Pair joint accuracy} To separate stable-and-correct behavior from stable-but-wrong behavior, the pair joint accuracy is defined by
\begin{equation}
\mathrm{PJA} =
\frac{1}{M}
\sum_{m=1}^{M}
\mathbb{I}[\hat{y}_{i_m}=y_{i_m} \;\land\; \hat{y}_{j_m}=y_{j_m}],
\label{eq:pja}
\end{equation}
which requires both members of the pair to be classified correctly.

\noindent\textbf{Mean normalized score drift} Because a detector can keep the same hard decision while still showing substantial score-level sensitivity, we also report the mean normalized score drift. Let
\begin{equation}
S = \{s_1,\ldots,s_N\},
\qquad
\mathrm{IQR}(S)=Q_{0.75}(S)-Q_{0.25}(S).
\label{eq:iqr_definition}
\end{equation}
Then
\begin{equation}
\mathrm{MNSD} =
\frac{1}{\max(\mathrm{IQR}(S), 10^{-12})}
\cdot
\frac{1}{M}
\sum_{m=1}^{M}
|s_{i_m}-s_{j_m}|.
\label{eq:mnsd}
\end{equation}
In the implementation, if the score IQR is numerically zero, the denominator is replaced by $1$; Eq.~\eqref{eq:mnsd} writes the same safeguard in compact form.

\noindent\textbf{Symmetric misclassification risk} The scorer also records an auxiliary diagnostic called symmetric misclassification risk. Define the sample-level 0--1 loss as
\begin{equation}
\ell_i = \mathbb{I}[\hat{y}_i \neq y_i].
\label{eq:zero_one_loss}
\end{equation}
The corresponding pairwise quantity is
\begin{equation}
\mathrm{SMR} =
\frac{1}{M}
\sum_{m=1}^{M}
\frac{\ell_{i_m}+\ell_{j_m}}{2}.
\label{eq:smr}
\end{equation}
SMR summarizes the average misclassification burden within a matched pair and is mainly useful as a diagnostic complement to PCR, PJA, and MNSD.

\subsection{Delivery-lineage metrics}

A delivery lineage groups all rendered samples that share the same parent utterance, delivery family, and binary label. Each node $v$ in lineage $g$ carries an observed operator-signature sequence $\sigma(v)$. The reference node $r_g$ is the shortest observed signature realization in that lineage.

\noindent\textbf{Lineage graph and node correctness} Distances are defined on a lineage graph rather than directly on template identifiers. The graph nodes are observed operator-signature realizations, and an undirected edge is inserted whenever two realizations differ by exactly one atomic delivery edit: a single parameter perturbation, a single operator substitution, an adjacent order swap, or a single insertion/deletion. Let the shortest-path distance from $r_g$ to node $v$ be $d_g(v)$.
Using the same fixed reference threshold from Eq.~\eqref{eq:pair_threshold}, node correctness is defined by
\begin{equation}
c(v) = \mathbb{I}[\hat{y}(v)=y(v)].
\label{eq:node_correctness}
\end{equation}

\noindent\textbf{Censored first-failure distance} For lineage $g$, the first-failure distance is
\begin{equation}
\rho_g = \min\{d_g(v): c(v)=0\},
\label{eq:ffd}
\end{equation}
whenever at least one misclassified node is observed in that lineage. Some lineages may have no observed failure at all and are therefore right-censored.

Let
\begin{equation}
D_{\max} = \max_{g \in \mathcal{G}}\; \max_{v \in g} d_g(v)
\label{eq:dmax}
\end{equation}
be the largest observed lineage distance in the evaluation set. The censored first-failure distance is then
\begin{equation}
\tilde{\rho}_g =
\begin{cases}
\rho_g, & \text{if lineage } g \text{ contains at least one observed failure},\\
D_{\max}+1, & \text{otherwise},
\end{cases}
\label{eq:censored_ffd_lineage}
\end{equation}
and the reported aggregate is
\begin{equation}
\mathrm{C\mbox{-}FFD} =
\frac{1}{|\mathcal{G}|}
\sum_{g \in \mathcal{G}}
\tilde{\rho}_g.
\label{eq:cffd}
\end{equation}
C-FFD increases when correct behavior persists farther from the reference node along the lineage graph.

\noindent\textbf{Robust-at-distance} To summarize robustness at progressively larger edit distances, we define
\begin{equation}
R_k =
\frac{1}{|\mathcal{G}|}
\sum_{g \in \mathcal{G}}
\mathbb{I}\!\left[
\forall v \in g,\ d_g(v) \le k \Rightarrow c(v)=1
\right],
\qquad k=0,1,\ldots,D_{\max}.
\label{eq:robust_at_k}
\end{equation}
Hence, $R_k$ is the fraction of lineages that remain entirely correct up to and including graph distance $k$.

\noindent\textbf{AURC-chain} Finally, the area under the robustness curve is defined by
\begin{equation}
\mathrm{AURC\mbox{-}chain} =
\frac{1}{D_{\max}+1}
\sum_{k=0}^{D_{\max}} R_k.
\label{eq:aurc_chain}
\end{equation}
AURC-chain compresses the full robustness profile into a single scalar and increases when correctness is maintained more consistently under progressively larger delivery edits.

% \appendix

% \section{Research Methods}

% \subsection{Part One}

% Lorem ipsum dolor sit amet, consectetur adipiscing elit. Morbi
% malesuada, quam in pulvinar varius, metus nunc fermentum urna, id
% sollicitudin purus odio sit amet enim. Aliquam ullamcorper eu ipsum
% vel mollis. Curabitur quis dictum nisl. Phasellus vel semper risus, et
% lacinia dolor. Integer ultricies commodo sem nec semper.

% \subsection{Part Two}

% Etiam commodo feugiat nisl pulvinar pellentesque. Etiam auctor sodales
% ligula, non varius nibh pulvinar semper. Suspendisse nec lectus non
% ipsum convallis congue hendrerit vitae sapien. Donec at laoreet
% eros. Vivamus non purus placerat, scelerisque diam eu, cursus
% ante. Etiam aliquam tortor auctor efficitur mattis.

% \section{Online Resources}

% Nam id fermentum dui. Suspendisse sagittis tortor a nulla mollis, in
% pulvinar ex pretium. Sed interdum orci quis metus euismod, et sagittis
% enim maximus. Vestibulum gravida massa ut felis suscipit
% congue. Quisque mattis elit a risus ultrices commodo venenatis eget
% dui. Etiam sagittis eleifend elementum.

% Nam interdum magna at lectus dignissim, ac dignissim lorem
% rhoncus. Maecenas eu arcu ac neque placerat aliquam. Nunc pulvinar
% massa et mattis lacinia.

\end{document}
\endinput
%%
%% End of file `sample-sigconf-authordraft.tex'.
